\section{Introducci\'on}
\label{sec:introduction}

\subsection{Motivaci\'on}

Cuando una muestra de malware llega a un equipo de respuesta a
incidentes o de inteligencia de amenazas, el resultado \'util rara vez
es el propio binario: es un resumen estructurado y accionable de qu\'e
indicadores de compromiso (IOCs) presenta, qu\'e t\'ecnicas de
MITRE ATT\&CK~\cite{mitre-attack} implementa, y artefactos listos para
alimentar ingenier\'ia de detecci\'on (reglas Sigma, firmas YARA) o una
plataforma de inteligencia de amenazas como MISP~\cite{misp}. A menudo
ya existe una primera pasada r\'apida sobre ese resumen: una consulta
de hash a un agregador de reputaci\'on como VirusTotal devuelve
veredictos de proveedores y etiquetas de comportamiento de sandbox en
segundos. Lo que esa primera pasada no ofrece es una respuesta
\emph{reconciliada}, ya que las muchas se\~nales que muestra se
presentan una junto a otra en lugar de contrastarse entre s\'i, de
modo que convertirlas en un mapeo de t\'ecnica defendible para una
muestra concreta todav\'ia suele significar que un analista humano
detona la muestra manualmente en un sandbox, traza su c\'odigo en un
desensamblador y redacta el informe a mano. Ese proceso es riguroso
pero lento, y no escala al volumen de muestras que un equipo de
operaciones de seguridad encuentra realmente.

Los sandboxes automatizados como CAPE~\cite{cape} ya producen grandes
vol\'umenes de datos de comportamiento en bruto: trazas de llamadas a
API, ficheros soltados, actividad de red, coincidencias de firmas
mantenidas por la comunidad. Las herramientas de an\'alisis est\'atico
extraen importaciones, cadenas y artefactos de cargas empaquetadas
directamente del binario. Lo que generalmente falta es la capa que se
sit\'ua entre estos hallazgos en bruto y algo que un analista o una
plataforma consumidora pueda usar directamente: reconciliar evidencia
est\'atica y din\'amica para la \emph{misma} afirmaci\'on, mapear esa
evidencia sobre una taxonom\'ia de t\'ecnicas est\'andar con un nivel de
confianza defendible, y exportarla en un formato que las plataformas de
inteligencia de amenazas ya consumen.

\subsection{Planteamiento del problema}

Tres problemas concretos motivan este trabajo, cada uno identificado al
construir y validar la tuber\'ia descrita en esta tesis:

\begin{enumerate}
  \item \textbf{Reconciliaci\'on de confianza entre fuentes.} Una
    t\'ecnica observada tanto por an\'alisis est\'atico (una
    combinaci\'on de importaciones) como por an\'alisis din\'amico (una
    firma de sandbox) es evidencia m\'as s\'olida que cualquiera de las
    dos por separado, pero solo si ambas observaciones se refieren
    realmente a la misma t\'ecnica. Una implementaci\'on ingenua que
    compara observaciones por identificador de t\'ecnica exacto pasa
    por alto el caso habitual en el que una fuente reporta una t\'ecnica
    padre y la otra reporta una subt\'ecnica espec\'ifica para lo que,
    en realidad, es el mismo comportamiento.
  \item \textbf{Malware multietapa.} Un cargador que suelta un driver
    rootkit, un m\'odulo antiantivirus y un cliente de acceso remoto
    tiene capacidades reales distribuidas entre varios binarios, no
    concentradas en el \'unico fichero detonado inicialmente. Fusionar
    ingenuamente la evidencia de cada componente soltado en la
    puntuaci\'on de confianza propia de la muestra padre atribuir\'ia
    err\'oneamente esa capacidad, ya que el comportamiento de la carga
    soltada no demuestra que el cargador realice ese comportamiento.
  \item \textbf{Confiar en el resultado.} Una tuber\'ia que reporta alta
    confianza para un hallazgo equivocado es peor que una que no
    reporta nada, porque enga\~na activamente a quien consume el
    resultado. Esto motiva tratar cada regla de detecci\'on como una
    afirmaci\'on que necesita evidencia, y cada resultado de
    evaluaci\'on como algo que auditar antes de confiar en \'el, no
    solo como una cifra que reportar.
\end{enumerate}

\subsection{Contribuciones}

Esta tesis aporta las siguientes contribuciones, cada una sustanciada
en detalle en secciones posteriores:

\begin{itemize}
  \item Un \textbf{modelo de reconciliaci\'on de confianza cruzada entre
    fuentes y entre niveles de t\'ecnica} (\S\ref{sec:reconciliation})
    que extiende la regla est\'andar ``ambas fuentes coinciden
    $\Rightarrow$ confianza alta'' para reconocer tambi\'en la
    concordancia entre una t\'ecnica padre reportada por una fuente y su
    subt\'ecnica espec\'ifica reportada por la otra, un hueco confirmado
    en datos reales de evaluaci\'on, no meramente te\'orico.
  \item Un \textbf{bucle de reenv\'io} (\S\ref{sec:resubmission}) que
    analiza de forma independiente cada fichero que una muestra
    multietapa suelta o inyecta durante la detonaci\'on, etiquetando
    cada uno con su linaje hacia la muestra padre sin mezclar su
    evidencia en la salida de confianza propia de la muestra padre.
  \item Un \textbf{conjunto de reglas de detecci\'on est\'atica
    reproducible y trazable a evidencia} (\S\ref{sec:coverage}),
    ampliado sistem\'aticamente de unas 30 a aproximadamente 85
    identificadores de t\'ecnica MITRE ATT\&CK relevantes para Windows,
    con cada adici\'on trazada a la descripci\'on oficial de MITRE para
    esa t\'ecnica y verificada frente al conjunto de muestras validadas
    antes de confiar en que no sobreajusta.
  \item Una \textbf{metodolog\'ia de evaluaci\'on y proceso de
    auditor\'ia} (\S\ref{sec:evaluation}) que va m\'as all\'a de
    precisi\'on/exhaustividad: precisi\'on estratificada por nivel de
    confianza y por acuerdo entre fuentes para comprobar si las propias
    afirmaciones del modelo de reconciliaci\'on se sostienen, y una
    auditor\'ia sistem\'atica de falsos positivos en la que cada
    discrepancia entre la salida de la tuber\'ia y la verdad de
    referencia verificada manualmente se diagnostic\'o individualmente a
    una causa concreta y verificable antes de corregirse, dejarse
    abierta, o atribuirse a un hueco en la verdad de referencia.
\end{itemize}

\subsection{Estructura del documento}

La Secci\'on~\ref{sec:related-work} sit\'ua este trabajo respecto a las
herramientas de sandboxing existentes, los enfoques de mapeo est\'atico
a ATT\&CK, y agregadores de reputaci\'on como VirusTotal. La
Secci\'on~\ref{sec:architecture} describe la arquitectura de la
tuber\'ia de principio a fin. La Secci\'on~\ref{sec:evaluation}
presenta la metodolog\'ia de evaluaci\'on y los resultados, incluida la
auditor\'ia de falsos positivos. La Secci\'on~\ref{sec:case-studies}
detalla cuatro casos concretos, incluyendo un error arquitect\'onico real
encontrado y corregido a trav\'es del proceso de evaluaci\'on, que
ilustran la disciplina de validaci\'on del proyecto con m\'as
concreci\'on de lo que permiten las cifras agregadas. La
Secci\'on~\ref{sec:limitations} enuncia con precisi\'on las
limitaciones de la tuber\'ia, cada una ligada a una causa espec\'ifica
diagnosticada. La Secci\'on~\ref{sec:future-work} esboza extensiones
concretas, incluyendo una, an\'alisis din\'amico Linux/ELF para
familias de botnets IoT, evaluada y deliberadamente aplazada por
razones expl\'icitas, no impl\'icitas. La Secci\'on~\ref{sec:conclusion}
concluye. El Anexo~\ref{app:rules} tabula el conjunto de reglas de
detecci\'on a\~nadidas durante la ampliaci\'on de cobertura. El
c\'odigo fuente completo de la tuber\'ia se publica en el repositorio
p\'ublico del proyecto en lugar de reproducirse aqu\'i.
